%----------------------------------------------------------------------------------------
% DE BINOMIALE VERDELING. Alle opgaven zijn eigen werk.
%
% Opbouw: Start, leerstof (theorie, werkwijze of uitleg, voorbeeld, korte oefening), en op een
% nieuwe pagina de Opdrachten (door elkaar, stijgende moeilijkheid), Koppel en Test jezelf.
% Labels van dit hoofdstuk beginnen met h4-.
%----------------------------------------------------------------------------------------

\mychapter{4}{De binomiale verdeling}

\lkr{Leerplan B+S: LPD 29 (kansen berekenen en interpreteren met de binomiale verdeling), met de wenken over ICT, verwachtingswaarde en standaardafwijking, en het grafische verband met de normale verdeling. Die normale verdeling zelf (kansen berekenen) hoort bij de cursus Statistiek; hier blijft het bij de figuur. Op de TI-84: \texttt{binompdf(n,p,k)} en \texttt{binomcdf(n,p,k)} in het menu DISTR.}

Een leerling gokt op een meerkeuzetoets. Een speler neemt vijf strafschoppen. Een fabriek controleert vijftig lampen. Telkens herhaal je hetzelfde experiment een aantal keer, met elke keer dezelfde kans op succes, en tel je de successen. Het aantal successen volgt dan altijd hetzelfde patroon: de binomiale verdeling. In dit hoofdstuk leer je de formule, het gebruik van je rekentoestel, en hoe je beoordeelt of een uitkomst verrassend is.

\begin{oefeningen}[Start]
\oef{h4-start-1}

\oef{h4-verken}
\end{oefeningen}

\lkr{Laat oefening \ref{oef:h4-verken} maken vóór de uitleg. Leerlingen tellen de paden met precies twee successen in een boomdiagram, en zien dat elk pad dezelfde kans heeft: de formule is dan \textit{aantal paden maal kans per pad}.}

%========================================================================================
\section{De formule}

\subsection{Binomiale experimenten}

Wat hebben een gokkende leerling, een strafschoppenreeks en een lampencontrole gemeen? Telkens herhaal je hetzelfde experiment, met maar twee uitkomsten die je interesseren, en tel je hoe vaak de ene uitkomst voorkomt. Zo'n situatie krijgt een naam.

\begin{theorie}
Een \textbf{\textcolor{accent}{binomiaal experiment}} bestaat uit $n$ herhalingen van een kansexperiment met twee uitkomsten, \textit{succes} en \textit{mislukking}, waarbij

\vspace{1mm}
\quad 1.\ \ de kans op succes $p$ bij elke herhaling dezelfde is, en

\quad 2.\ \ de herhalingen onafhankelijk zijn.

\vspace{1mm}
Het aantal successen $X$ heeft dan een \textbf{\textcolor{accent}{binomiale verdeling}} met parameters $n$ en $p$, en je schrijft $\formule{X \sim B(n, p)}$.
\end{theorie}

\textit{Succes} is gewoon de uitkomst die je telt, ook als het iets slechts is, zoals een defecte lamp. Voor je rekent, ga je na of beide voorwaarden kloppen. De tweede valt het vaakst weg: trek je zonder teruglegging uit een kleine groep, dan hangt elke trekking af van de vorige. Uit een grote groep mag het wel bij benadering: of één lamp van 10\,000 eruit is, verandert de kans op een defecte lamp bijna niet.

\voorbeeld{}{Is $X$ binomiaal? Geef dan $n$ en $p$. (a) $X$ is het aantal keer kop bij 15 worpen met een munt. (b) Uit een vaas met 3 rode en 2 witte knikkers trek je er 2 zonder teruglegging; $X$ is het aantal rode. (c) $X$ is het aantal juiste antwoorden van een leerling op een toets van 10 vragen, die steeds moeilijker worden.\label{vb:h4-model}}{%
\uitleg{We zoeken telkens een vast aantal herhalingen met dezelfde kans, die elkaar niet beïnvloeden.}
\begin{stappenvrij}{3.4cm}{11.2cm}
\text{(a) } X \sim B(15;\, 0{,}5) & Vijftien onafhankelijke worpen, telkens kans $\frac12$ op kop. \\[1ex]
\text{(b) niet binomiaal} & Is de eerste rood, dan is de tweede rood met kans $\frac24$, anders met kans $\frac34$: de trekkingen zijn afhankelijk. Zo is $P(X = 2) = \frac35 \cdot \frac24 = \frac{3}{10}$, terwijl de formule met $n = 2$ en $p = \frac35$ ten onrechte $\frac{9}{25}$ zou geven. \\[1ex]
\text{(c) niet binomiaal} & De kans op een juist antwoord is niet bij elke vraag dezelfde.
\end{stappenvrij}}

\begin{oefeningen}
\oef{h4-model}
\end{oefeningen}

\needspace{16\baselineskip}
\subsection{De binomiale kans}

In de Verken vond je de kans op precies twee keer scoren bij vier vrijworpen: elk rijtje met twee keer scoren heeft dezelfde kans $0{,}7^2 \cdot 0{,}3^2$, en er zijn $\binom42 = 6$ zulke rijtjes. Die redenering werkt voor elke $n$, $p$ en $k$.

\begin{theorie}
\begin{stelling}{Binomiale kans}
Voor $X \sim B(n, p)$ en $k = 0, 1, \ldots, n$ is \quad $\formule{P(X = k) = \binom{n}{k} p^k (1 - p)^{n - k}}$
\end{stelling}

\begin{bewijs}
Schrijf de uitkomst van de $n$ herhalingen als een rijtje van $S$ en $M$. Neem een rijtje met precies $k$ keer $S$, zoals $SS\!M\!S\!M$. Omdat de herhalingen onafhankelijk zijn, is zijn kans het product van $k$ factoren $p$ en $n - k$ factoren $1 - p$:
\bewijsregel{P(\text{één rijtje met } k \text{ successen}) = p^k (1 - p)^{n - k} .}
Zo'n rijtje ligt vast als je kiest op welke $k$ van de $n$ plaatsen een $S$ staat. Dat kan op $\binom{n}{k}$ manieren, dus er zijn $\binom{n}{k}$ zulke rijtjes, allemaal met dezelfde kans. Twee verschillende rijtjes zijn onverenigbaar, dus de somregel geeft
\bewijsregel{P(X = k) = \binom{n}{k} \cdot p^k (1 - p)^{n - k} .}
\end{bewijs}
\end{theorie}

In woorden: \textit{aantal rijtjes maal de kans op één rijtje}. De factor $p^k(1 - p)^{n - k}$ zegt hoe waarschijnlijk één bepaalde volgorde is, de factor $\binom{n}{k}$ telt de volgordes.

\voorbeeld{}{Een leerling gokt op 10 meerkeuzevragen met elk 4 antwoorden. Wat is de kans op precies 5 juiste antwoorden?\label{vb:h4-formule}}{%
\uitleg{Elke vraag is een herhaling met succeskans $\frac14$, onafhankelijk van de andere. Het aantal juiste antwoorden is binomiaal met $n = 10$ en $p = 0{,}25$.}
\begin{stappenbreed}
X \sim B(10;\, 0{,}25) & De herhalingen zijn onafhankelijk, met dezelfde kans. \\[1ex]
P(X = 5) = \binom{10}{5} \cdot 0{,}25^5 \cdot 0{,}75^5 = 252 \cdot 0{,}25^5 \cdot 0{,}75^5 \approx 0{,}058 & 252 volgordes van 5 juiste en 5 foute. Op het rekentoestel: \texttt{binompdf(10,0.25,5)}.
\end{stappenbreed}}

\nujij{Je gooit 8 keer met een munt. Wat is de kans op precies 4 keer kop?}{$\binom84 \cdot 0{,}5^8 = \frac{70}{256} \approx 0{,}273$.}

\begin{oefeningen}
\oef{h4-formule-2}
\end{oefeningen}

\needspace{16\baselineskip}
\subsection{Cumulatieve kansen}

Vraag je naar \textit{hoogstens} of \textit{minstens}, dan tel je binomiale kansen op. Met de hand is dat veel werk. Je rekentoestel geeft de cumulatieve verdelingsfunctie $F(k) = P(X \leqslant k)$ meteen, met \texttt{binomcdf(n,p,k)}. De rest volgt met de regels uit hoofdstuk 3.

\voorbeeld{}{Wat is de kans dat de leerling uit het vorige voorbeeld minstens 6 juiste antwoorden gokt, en dus slaagt?\label{vb:h4-cumulatief}}{%
\uitleg{\textit{Minstens 6} is het complement van \textit{hoogstens 5}, en dat geeft het rekentoestel als cumulatieve kans.}
\begin{stappenbreed}
P(X \geqslant 6) = 1 - P(X \leqslant 5) & Het complement. \\[1ex]
= 1 - \texttt{binomcdf(10,0.25,5)} \approx 1 - 0{,}980 = 0{,}020 & Ongeveer 2~\%: slagen door te gokken is onwaarschijnlijk.
\end{stappenbreed}}

\voorbeeld{}{Wat is de kans dat dezelfde leerling 3, 4 of 5 juiste antwoorden gokt?\label{vb:h4-tussen}}{%
\uitleg{Een kans tussen twee grenzen is een verschil van twee cumulatieve kansen. Let op welke waarde je aftrekt: 3 moet erbij blijven.}
\begin{stappenbreed}
P(3 \leqslant X \leqslant 5) = P(X \leqslant 5) - P(X \leqslant 2) & Van \textit{hoogstens 5} trek je \textit{hoogstens 2} af. \\[1ex]
\approx 0{,}980 - 0{,}526 = 0{,}455 & \texttt{binomcdf(10,0.25,5)} min \texttt{binomcdf(10,0.25,2)}.
\end{stappenbreed}
Wie $P(X \leqslant 3)$ aftrekt, vergeet $P(X = 3) \approx 0{,}250$, de grootste van de drie kansen.}

\begin{werkwijze}
Ga na: zijn er $n$ onafhankelijke herhalingen met telkens dezelfde kans $p$ op succes? Bepaal $n$, $p$ en de gevraagde $k$. Voor $P(X = k)$: de formule of \texttt{binompdf}. Voor $P(X \leqslant k)$: \texttt{binomcdf}. Voor $P(X \geqslant k)$: $1 - P(X \leqslant k - 1)$. Voor $P(k \leqslant X \leqslant m)$: $P(X \leqslant m) - P(X \leqslant k - 1)$. Let op de grens: \textit{minder dan 4} is $X \leqslant 3$.
\end{werkwijze}

\begin{oefeningen}
\oef{h4-formule}

\oef{h4-grens}
\end{oefeningen}

%========================================================================================
\section{Verwachtingswaarde en standaardafwijking}

\subsection{Een rekenregel voor $\binom{n}{k}$}

Gokt de leerling op 10 vragen met kans $\frac14$, dan verwacht je gemiddeld $10 \cdot \frac14 = 2{,}5$ juiste antwoorden. Om te bewijzen dat dat altijd zo werkt, hebben we één eigenschap van $\binom{n}{k}$ nodig.

\begin{theorie}
\begin{stelling}{Kapiteinsregel}
Voor $1 \leqslant k \leqslant n$ is \quad $\formule{k \binom{n}{k} = n \binom{n - 1}{k - 1}}$
\end{stelling}

\begin{bewijs}
Tel op twee manieren hoeveel ploegen van $k$ spelers met één kapitein je kunt vormen uit $n$ personen. Kies je eerst de ploeg en dan de kapitein in de ploeg, dan zijn er $\binom{n}{k}$ ploegen, met elk $k$ mogelijke kapiteins: samen $k\binom{n}{k}$. Kies je eerst de kapitein uit de $n$ personen en dan de $k - 1$ andere spelers uit de $n - 1$ overigen, dan vind je $n \binom{n - 1}{k - 1}$. Beide tellen dezelfde ploegen, dus de aantallen zijn gelijk.
\end{bewijs}
\end{theorie}

Ter controle met $n = 5$ en $k = 2$: $2 \cdot \binom52 = 20$ en $5 \cdot \binom41 = 20$. Met de formules hieronder vind je voor het aantal zessen bij twee worpen ($n = 2$, $p = \frac16$) dezelfde $\mu = \frac13$ en $\Var(X) = \frac{5}{18}$ als in hoofdstuk 3.

\needspace{16\baselineskip}
\subsection{De formules}

\begin{theorie}
\begin{stelling}{Verwachtingswaarde van de binomiale verdeling}
Voor $X \sim B(n, p)$ is \quad $\formule{\E(X) = np}$
\end{stelling}

\begin{bewijs}
Schrijf $q = 1 - p$. De term met $k = 0$ is nul, en voor de andere termen geeft de kapiteinsregel $k\binom{n}{k} = n\binom{n-1}{k-1}$:
\bewijsregel{\E(X) = \sum_{k=1}^{n} k \binom{n}{k} p^k q^{n-k} = \sum_{k=1}^{n} n \binom{n-1}{k-1} p^k q^{n-k} = np \sum_{k=1}^{n} \binom{n-1}{k-1} p^{k-1} q^{(n-1)-(k-1)} .}
Stel $j = k - 1$. Dan loopt $j$ van 0 tot $n - 1$, en is de laatste som $\sum_{j=0}^{n-1} \binom{n-1}{j} p^j q^{n-1-j}$. Dat zijn alle kansen van $B(n - 1, p)$, en die zijn samen 1. Dus $\E(X) = np$.
\end{bewijs}

\begin{stelling}{Standaardafwijking van de binomiale verdeling}
Voor $X \sim B(n, p)$ is \quad $\formule{\Var(X) = np(1 - p)}$ \quad en \quad $\formule{\sigma = \sqrt{np(1 - p)}}$
\end{stelling}

\begin{bewijs}
Voor $n = 1$ neemt $X$ alleen de waarden 0 en 1 aan, dus $\E(X^2) = \E(X) = p$ en $\Var(X) = p - p^2 = p(1 - p)$. Neem nu $n \geqslant 2$. De kapiteinsregel twee keer toegepast geeft, voor $k \geqslant 2$,
\bewijsregel{k(k - 1)\binom{n}{k} = (k - 1) \cdot n\binom{n - 1}{k - 1} = n(n - 1)\binom{n - 2}{k - 2} .}
Net als hierboven is dan, met $j = k - 2$ en de kansen van $B(n - 2, p)$ die samen 1 zijn,
\bewijsregel{\sum_{k=0}^{n} k(k - 1)\, P(X = k) = n(n - 1)p^2 \sum_{j=0}^{n-2} \binom{n-2}{j} p^j q^{n-2-j} = n(n - 1)p^2 .}
Omdat $k^2 = k(k - 1) + k$, is $\E(X^2) = n(n - 1)p^2 + np$. De rekenformule voor de variantie geeft
\bewijsregel{\Var(X) = n(n - 1)p^2 + np - (np)^2 = np - np^2 = np(1 - p) .}
\end{bewijs}
\end{theorie}

De eerste formule is wat je intuïtief verwacht: van $n$ pogingen lukt op lange termijn een deel $p$. De spreiding is het grootst bij $p = \frac12$, en klein als $p$ dicht bij 0 of 1 ligt. Ze groeit maar met $\sqrt{n}$: het aandeel successen $\frac{X}{n}$ schommelt dus minder naarmate $n$ groter wordt, zoals de wet van de grote aantallen zegt.

\voorbeeld{}{Een zaadje kiemt met kans 0,9. Je zaait er 12. Hoeveel kiemen er gemiddeld, met welke standaardafwijking? Wat is de kans dat ze allemaal kiemen?\label{vb:h4-mu}}{%
\uitleg{De zaadjes kiemen onafhankelijk van elkaar met dezelfde kans: $X \sim B(12;\, 0{,}9)$. Dan gebruiken we de formules in plaats van de definitie.}
\begin{stappen}
\mu = 12 \cdot 0{,}9 = 10{,}8 & Gemiddeld bijna 11 van de 12. \\[1ex]
\sigma = \sqrt{12 \cdot 0{,}9 \cdot 0{,}1} \approx 1{,}04 & Kleine spreiding, want $p$ ligt dicht bij 1. \\[1ex]
P(X = 12) = 0{,}9^{12} \approx 0{,}28 & Alle twaalf: maar in 28~\% van de gevallen.
\end{stappen}}

\begin{oefeningen}
\oef{h4-mu}

\oef{h4-mu-2}
\end{oefeningen}

%========================================================================================
\section{Vorm van de verdeling}

\subsection{De normale verdeling als benadering}

Hoe groter $n$, hoe meer het staafdiagram van een binomiale verdeling op een klok lijkt, ook als $p$ niet $\frac12$ is. Bij $B(5;\, 0{,}3)$ is het staafdiagram nog duidelijk scheef, bij $B(15;\, 0{,}3)$ al veel minder.

\figuren{\begin{tikzpicture}[x=.62cm,y=6cm,font=\scriptsize]
  \draw[black!14,very thin] (-.5,0) grid[xstep=20,ystep=.1] (5.5,.4);
  \foreach \k/\h in {0/.1681,1/.3601,2/.3087,3/.1323,4/.0283,5/.0024}{\staaf{\k}{\h}{.6}}
  \draw[fel!80!black,thick] plot[domain=-.5:5.5,samples=60,smooth] (\x,{exp(-((\x-1.5)^2)/2.1)/sqrt(2*pi*1.05)});
  \draw[->,grijs] (-.5,0) -- (5.8,0) node[right,text=black] {$k$};
  \foreach \k in {0,...,5}{\node[below] at (\k,0) {\k};}
  \node[anchor=north east] at (5.5,.42) {$B(5;\, 0{,}3)$};
\end{tikzpicture}\qquad
\begin{tikzpicture}[x=.36cm,y=6cm,font=\scriptsize]
  \draw[black!14,very thin] (-.5,0) grid[xstep=20,ystep=.1] (12.5,.4);
  \foreach \k/\h in {0/.0047,1/.0305,2/.0916,3/.1700,4/.2186,5/.2061,6/.1472,7/.0811,8/.0348,9/.0116,10/.0030,11/.0006}{\staaf{\k}{\h}{.6}}
  \draw[fel!80!black,thick] plot[domain=-.5:12.5,samples=80,smooth] (\x,{exp(-((\x-4.5)^2)/6.3)/sqrt(2*pi*3.15)});
  \draw[->,grijs] (-.5,0) -- (12.8,0) node[right,text=black] {$k$};
  \foreach \k in {0,2,...,12}{\node[below] at (\k,0) {\k};}
  \node[anchor=north east] at (12.5,.42) {$B(15;\, 0{,}3)$};
\end{tikzpicture}}

De klokvorm is de grafiek van de \textbf{normale verdeling} met hetzelfde gemiddelde en dezelfde standaardafwijking. De figuur hieronder toont $B(40;\, 0{,}3)$, met $\mu = 12$ en $\sigma \approx 2{,}9$, en de normale kromme: die volgt de staven bijna perfect.

\figuren{\begin{tikzpicture}[x=.48cm,y=24cm,font=\scriptsize]
  \draw[black!14,very thin] (0,0) grid[xstep=40,ystep=.05] (26,.15);
  \foreach \k/\h in {2/.0001,3/.0005,4/.002,5/.0061,6/.0151,7/.0315,8/.0557,9/.0849,10/.1128,11/.1319,12/.1366,13/.1261,14/.1042,15/.0774,16/.0518,17/.0314,18/.0172,19/.0085,20/.0038,21/.0016,22/.0006}{\staaf{\k}{\h}{.7}}
  \draw[fel!80!black,very thick] plot[domain=1:25,samples=120,smooth] (\x,{exp(-((\x-12)^2)/16.8)/sqrt(2*pi*8.4)});
  \draw[->,grijs] (0,0) -- (26.5,0) node[right,text=black] {$k$};
  \draw[->,grijs] (0,0) -- (0,.16);
  \foreach \k in {0,4,...,24}{\node[below] at (\k,0) {\k};}
  \foreach \t/\l in {.05/{0,05},.1/{0,10},.15/{0,15}}{\node[left] at (0,\t) {\l};}
\end{tikzpicture}}

Als vuistregel is de benadering goed als $np \geqslant 5$ en $n(1 - p) \geqslant 5$. Bij $B(40;\, 0{,}3)$ is $np = 12$ en $n(1 - p) = 28$: goed. Bij $B(5;\, 0{,}3)$ is $np = 1{,}5$: de benadering is slecht, en dat zie je ook. Later leer je met de normale verdeling rekenen. Voor de binomiale verdeling zelf gebruik je gewoon je rekentoestel.

\needspace{16\baselineskip}
\subsection{Is een uitkomst verrassend?}

\naast{\begin{minipage}{7.6cm}
Past een uitkomst nog bij het toeval? Een eerste idee geeft het interval $\mu \pm 2\sigma$. Neem de gokkende leerling: $B(10;\, 0{,}25)$ heeft zijn top bij 2, dicht bij $\mu = 2{,}5$. Bijna alle kansmassa ligt binnen $\mu \pm 2\sigma$, met $\sigma = \sqrt{10 \cdot 0{,}25 \cdot 0{,}75} \approx 1{,}37$: tussen ongeveer $-0{,}2$ en $5{,}2$. Een uitkomst ver buiten dat interval, zoals 8 juiste antwoorden, is bij gokken zeer onwaarschijnlijk.
\end{minipage}}{\begin{tikzpicture}[x=.62cm,y=8.5cm,font=\scriptsize]
  \draw[black!14,very thin] (-.5,0) grid[xstep=20,ystep=.1] (10.5,.3);
  \foreach \k/\h in {0/.0563,1/.1877,2/.2816,3/.2503,4/.1460,5/.0584,6/.0162,7/.0031,8/.0004,9/0,10/0}{\staaf{\k}{\h}{.6}}
  \draw[->,grijs] (-.5,0) -- (10.8,0) node[right,text=black] {$k$};
  \draw[->,grijs] (-.5,0) -- (-.5,.32);
  \foreach \k in {0,2,...,10}{\node[below] at (\k,0) {\k};}
  \foreach \t/\l in {.1/{0,1},.2/{0,2},.3/{0,3}}{\node[left] at (-.5,\t) {\l};}
  \draw[fel!80!black,thick] (2.5,0) -- (2.5,.3) node[above,text=black] {$\mu$};
\end{tikzpicture}}

Nauwkeuriger is een kans, berekend met de binomiale verdeling. Kijk daarvoor niet naar de kans op die ene uitkomst: bij 100 worpen met een eerlijke munt is zelfs de meest waarschijnlijke uitkomst, 50 keer kop, maar kans $0{,}08$. Kijk naar de kans op een uitkomst die \textit{minstens zo extreem} is. Is die heel klein, dan twijfel je aan het model.

\voorbeeld{}{Een munt landt 62 keer op kop bij 100 worpen. Is dat verrassend voor een eerlijke munt?\label{vb:h4-verrassend}}{%
\uitleg{Voor een eerlijke munt is $X \sim B(100;\, 0{,}5)$, met $\mu = 50$ en $\sigma = 5$. 62 ligt boven $\mu$, dus we berekenen de kans op 62 of meer keer kop.}
\begin{stappenbreed}
P(X \geqslant 62) = 1 - P(X \leqslant 61) \approx 0{,}010 & Met \texttt{binomcdf(100,0.5,61)}. \\[1ex]
62 \text{ ligt } \tfrac{62 - 50}{5} = 2{,}4 \text{ standaardafwijkingen boven } \mu & Ver buiten $\mu \pm 2\sigma$.
\end{stappenbreed}
Bij een eerlijke munt gebeurt dit maar in ongeveer 1~\% van de reeksen. Dat is een sterke aanwijzing dat de munt niet eerlijk is, maar geen bewijs.}

Ligt de uitkomst onder $\mu$, dan is \textit{minstens zo extreem} de andere kant op: je berekent $P(X \leqslant k)$.

\voorbeeld{}{Een basketbalspeler beweert dat hij 80~\% van zijn vrijworpen scoort. Bij een test scoort hij er 11 van de 20. Geloof je hem nog?\label{vb:h4-verrassend-laag}}{%
\uitleg{Als zijn bewering klopt, is $X \sim B(20;\, 0{,}8)$, met $\mu = 16$. Hij scoorde minder, dus berekenen we de kans op 11 of minder.}
\begin{stappenbreed}
\sigma = \sqrt{20 \cdot 0{,}8 \cdot 0{,}2} \approx 1{,}79 & 11 ligt $\frac{16 - 11}{1{,}79} \approx 2{,}8$ standaardafwijkingen onder $\mu$. \\[1ex]
P(X \leqslant 11) \approx 0{,}010 & Met \texttt{binomcdf(20,0.8,11)}.
\end{stappenbreed}
Als hij echt 80~\% scoort, haalt hij zo'n slechte reeks maar in ongeveer 1 op de 100 tests. Dat is een sterke reden om aan zijn bewering te twijfelen.}

\begin{werkwijze}
Stel het model op dat je wilt toetsen, $X \sim B(n, p)$, en bereken $\mu$. Ligt de uitkomst $k$ boven $\mu$, bereken dan $P(X \geqslant k)$; ligt ze eronder, dan $P(X \leqslant k)$. Is die kans heel klein (bijvoorbeeld onder 5~\%), dan is de uitkomst verrassend en twijfel je aan het model. Een kleine kans is geen bewijs.
\end{werkwijze}

\begin{oefeningen}
\oef{h4-verrassend}

\oef{h4-verrassend-2}

\oef{h4-verrassend-laag}
\end{oefeningen}

%========================================================================================
\opdrachtenkop

\begin{oefeningen}
\needspace{10\baselineskip}
\oef{h4-o11}

\oef{h4-o1}

\needspace{10\baselineskip}
\oef{h4-o3}

\oef{h4-o13}

\oef{h4-o12}

\oef{h4-o2}

\oef{h4-o15}

\needspace{16\baselineskip}
\oef{h4-o14}

\oef{h4-o4}

\needspace{10\baselineskip}
\oef{h4-o16}

\oef{h4-o17}

\oef{h4-o5}

\oef{h4-o20}

\oef{h4-o21}

\oef{h4-o7}

\oef{h4-o18}

\oef{h4-o19}

\oef{h4-o6}

\oef{h4-o22}

\oef{h4-o8}

\oef{h4-o9}

\needspace{10\baselineskip}
\oef{h4-o23}

\oef{h4-o10}

\oef{h4-o24}

\oef{h4-o25}
\end{oefeningen}

%----------------------------------------------------------------------------------------
% KOPPEL: verbanden met andere leerstof
%----------------------------------------------------------------------------------------
\needspace{20\baselineskip}
\opdrachtsectie{Koppel}

\lkr{De oefening gebruikt logaritmen (Exponentiële en logaritmische functies).}

\begin{oefeningen}

\oef{h4-koppel-2}
\end{oefeningen}

%----------------------------------------------------------------------------------------
\opdrachtsectie{Test jezelf}

\begin{oefeningen}
\oef{h4-t1}

\oef{h4-t2}

\oef{h4-t3}

\oef{h4-t4}

\oef{h4-t5}

\oef{h4-t6}
\end{oefeningen}

\herhaalschema{%
  \fv{oef:h4-t1}{vb:h4-formule}{oef:h4-formule-2}%
  \fv{oef:h4-t2}{vb:h4-cumulatief}{oef:h4-formule}%
  \fv{oef:h4-t3}{vb:h4-mu}{oef:h4-mu}%
  \fv{oef:h4-t4}{vb:h4-verrassend-laag}{oef:h4-verrassend-laag}%
  \fv{oef:h4-t5}{vb:h4-cumulatief}{oef:h4-grens}%
  \fv{oef:h4-t6}{vb:h4-mu}{oef:h4-mu-2}}

\ifzelfstudie
\section*{Maak je eigen samenvatting}
Sluit dit boek en schrijf in je schrift: wanneer een experiment binomiaal is, de formule voor $P(X = k)$ en waarom ze klopt, hoe je \textit{hoogstens} en \textit{minstens} met het rekentoestel berekent, de verwachtingswaarde en standaardafwijking, het verband met de normale verdeling, en hoe je beoordeelt of een uitkomst verrassend is. Vergelijk daarna met de kaders.
\fi
